\subsection*{Intended Learning Outcomes - Lecture 04}

We want the students to be able to:
\begin{itemize}
\item explain what a Markov kernel is and how to represent them in the discrete and absolute continuous case.
\item give examples of Markov kernels.
\item explain how Markov kernels relate to deterministic maps and probability distributions.
\item explain transition probability spaces and null-sets.
\item explain and work with the 5 constructions of Markov kernels from others 
    (composition, product, marginal, conditional, push-forward) 
    and how to work with them in the discrete and absolute continuous case.
%\item Write down the definition of conditional independence for random fields 
%    and work with them in the discrete and continuous cases.
%\item Write down the separoid axioms for conditional independence of random fields.
%\item prove separoid axioms for random variables in the discrete and absolute continuous case (in $\R^D$).
%\item explain how one can get a given Markov kernel by a deterministic mapping and a probability distribution.
\end{itemize}
